\section{深圳多联机及风机盘管系统设计}

本章以深圳（夏热冬暖地区）为目标城市。深圳位于华南沿海，属于建筑热工设计气候分区中的夏热冬暖B区，全年温暖湿润，冬季短促温和，夏季漫长酷热。该城市在本文中代表华南沿海纯制冷导向型气候的极端样本。

\subsection{室外设计参数}

深圳夏季室外气象设计参数（CSWD气象年数据）为：夏季空调室外干球计算温度\SI{33.7}{\degreeCelsius}，夏季空调室外湿球计算温度\SI{27.5}{\degreeCelsius}，夏季空调室外计算日平均温度\SI{29.8}{\degreeCelsius}，夏季室外平均风速\SI{2.4}{\meter\per\second}。冬季空调室外干球计算温度为\SI{8.4}{\degreeCelsius}，为五城市中最高值，表明深圳的冬季供暖需求在五个城市中最为有限。深圳全年最热月平均温度（\SI{28.8}{\degreeCelsius}）为五城市中最高，全年制冷季极长，是本文代表“以制冷为全年主导”的典型城市。

\subsection{建筑概况}

同第一章详述。

\subsection{理想负荷计算结果}

\begin{table}[tbp]
\centering
\caption{深圳负荷计算与系统方案比较结果}
\label{tab:shenzhen-ideal}
\scriptsize
\begin{minipage}[t]{0.42\textwidth}
\centering
\textbf{（a）理想负荷工况结果}\\[3pt]
\begin{tabular}{@{}lr@{}}
\toprule
指标 & 数值 \\
\midrule
$Q_{\mathrm{pc}}$ (kW) & 15.74 \\
$q_{\mathrm{pc}}$ (W/m$^{2}$) & 11.51 \\
$E_{\mathrm{cool}}$ (kWh) & 27080 \\
$E_{\mathrm{heat}}$ (kWh) & 86479 \\
$E_{\mathrm{total}}$ (kWh) & 113559 \\
$E/A_{\mathrm{f}}$ (kWh/m$^{2}$) & 83.07 \\
\bottomrule
\end{tabular}
\end{minipage}%
\hfill%
\begin{minipage}[t]{0.54\textwidth}
\centering
\textbf{（b）两类系统比较}\\[3pt]
\begin{tabular}{@{}lcc@{}}
\toprule
指标 & VRF+DOAS & FCU+DOAS \\
\midrule
设计冷量 (kW) & 88.51 & 162.73 \\
设计热量 (kW) & 88.51 & 55.36 \\
年电耗 (kWh) & 76830 & 7746 \\
电耗强度 (kWh/m$^{2}$) & 56.20 & 5.67 \\
\bottomrule
\end{tabular}
\end{minipage}
\end{table}

深圳的峰值冷负荷为\SI{15.74}{\kilo\watt}，为五城市中最高之一；年冷负荷为\SI{27079.74}{\kilo\watt\hour}，同样居首。年冷负荷约为年热负荷的31.3\%（按统一围护结构计算，在当前差异化围护结构下，年冷负荷占总理想负荷的约23.8\%）。然而值得注意的是，即使在深圳这种以制冷为主的夏热冬暖地区，年热负荷仍达到\SI{86479.18}{\kilo\watt\hour}，占总负荷的约76.2\%。这主要源于两个因素：其一，建筑内部得热量在冬季仍需通过空调系统移除，热负荷中的相当部分实际上来自维持室内热舒适的“再热”需求；其二，由于深圳冬季无暖气，夜间低温时段新风和外窗得热的失热累积。这一结果提示，即使在夏热冬暖地区，建筑供热需求仍不可完全忽略。

\subsection{空调系统方案比较}

深圳VRF设计冷量为\SI{88.51}{\kilo\watt}，为五城市中最高，反映了该城市在全制冷工况下对设备容量的最大需求。VRF年电耗强度为\SI{56.20}{\kilo\watt\hour\per\square\meter}，低于沈阳（\SI{60.25}{\kilo\watt\hour\per\square\meter}）和天津（\SI{60.20}{\kilo\watt\hour\per\square\meter}）。这一看似反直觉的结果可作如下解读：沈阳和天津的VRF系统在漫长供暖季中以较低COP持续运行供热工况，全年压缩机累计运行小时数远超深圳以制冷为主的VRF系统；深圳VRF虽以全制冷模式运行为主，但制冷工况下VRF的COP通常远高于热泵工况，因此其单位冷量电耗更低。

深圳FCU冷水机组设计冷量（\SI{162.73}{\kilo\watt}）为五城市最高，锅炉设计热量（\SI{55.36}{\kilo\watt}）则接近成都水平。FCU冷量与热量之比为2.94:1，同样反映了显著的制冷主导特征。其年度天然气消耗为\SI{61653.32}{\kilo\watt\hour}（强度为\SI{45.10}{\kilo\watt\hour\per\square\meter}）。
